\section{Security}
\label{sec:security}

\label{sec:threat}

The participants are $\Nc$ clients and a server. The server is semi-honest. It
follows the protocol and may try to infer private information from what it
observes. Clients may likewise be semi-honest and may collude. \Cref{sec:malicious}
strengthens the client side of this model, and allows up to $\tc-1$ clients to
deviate arbitrarily while the server stays honest. \Cref{sec:malicious-ext}
explains why that assumption cannot be dropped.

\paragraph{View of the server} Ciphertexts under the collective public key, the
public per-client sample counts, and the evaluation keys. It holds no decryption
key. The server additionally sees
encrypted query features, which it cannot read.

The design is arranged so that no arithmetic shortcut evades it. No aggregate of
the adapters exists, so there is no sum from which a coalition can subtract its
own contributions. The shared head is never decrypted, so it cannot be inverted
the same way. The per-class totals are never decrypted, so a coalition cannot
recover the remaining client's class histogram.

\tronly{That last sentence is worth stating precisely, because it does not say the
histogram is unrecoverable. It says the protocol publishes no direct read of it.
A client's counts also shape the displacement it uploads, since for a linear map
trained under cross-entropy the magnitude of the row for a class grows with the
number of examples of that class the client held~\cite{wainakh2022user}. Those
displacements are ciphertexts and the merged head is a ciphertext, so nothing is
readable while the encryption stands, and \cref{sec:c2c-measured} measures what
survives into an extracted copy. What inverting the totals under encryption buys
is the removal of the one channel that would have handed the histogram over in
the clear.}

\tronly{This section states what the protocol guarantees. \Cref{sec:ideal} gives an
ideal functionality. \Cref{sec:semihonest} proves that the protocol realizes it
against semi-honest parties. \Cref{sec:malicious} bounds what a malicious client
coalition learns about an honest client's data. \Cref{sec:malicious-ext} shows
that the server must be honest and what would remove that requirement.
\Cref{sec:inherent} separates what the cryptography guarantees from what the
task itself leaks.}

\subsection{The Ideal Functionality}
\label{sec:ideal}

\Cref{func:ideal} defines the ideal functionality $\Fhe$ against which we
measure the protocol. The server is an entity outside $\Fhe$, so it may be corrupt. The form of the functionality follows the
secure aggregation functionality of ELSA~\cite{rathee2023elsa}, and placing the
aggregator outside it follows FULLSA~\cite{karakoc2024fullsa}.

\begin{functionality}[t]
\caption{HE-OFT $\Fhe$}
\label{func:ideal}
\footnotesize
\noindent\textbf{\emph{Parameters.}} The client count $\Nc$, the decryption
threshold $\tc$, the per-client query allowance $Q$, and the label space size
$\Cc$. The parties are clients $P_1,\dots,P_{\Nc}$ and a server $S$, which is not
part of $\Fhe$. The functionality holds a
counter $q_j$ for each client, each initialised to $0$, and a store that is
initially empty and that $\Fhe$ never sends to any party.

\noindent\textbf{\emph{Inputs.}} Each client $P_j$ inputs its dataset $\Dj$, and
later inputs queries $x$. The server $S$ inputs a selection request.

\noindent\textbf{\emph{Outputs.}} Every party receives the selected index
$a^{\star}$ and the phase signals. Client $P_j$ receives one label for each of
its first $Q$ queries, and $\bot$ afterwards.

\noindent\textbf{\emph{Steps.}}
\begin{enumerate}\itemsep2pt \parskip0pt \topsep3pt
\item \emph{Training and aggregation.} On $(\mathsf{Train},\Dj)$ from every
  client, run the local training of \cref{sec:split} on each $\Dj$ in both
  servable arrangements and store the two shared heads $\thstar_{A}$ and
  $\thstar_{B}$ of \cref{eq:aggregate}. Send
  $(\mathsf{Done},\mathsf{train})$ to every party.
\item \emph{Selection.} On $(\mathsf{Select})$ from $S$, evaluate the estimator
  of \cref{eq:estimator} on each of the two stored heads, send the index
  $a^{\star}\in\{A,B\}$ of the larger to every party, and send
  $(\mathsf{Done},\mathsf{select})$ to every party.
\item \emph{Serving.} On $(\mathsf{Query},x)$ from client $P_j$, return $\bot$
  to $P_j$ if $q_j\ge Q$. Otherwise set $q_j\leftarrow q_j+1$, return
  $\hat y=\arg\max_c\,(\thstar_{a^{\star}}\,\featx^{a^{\star}}_j(x))_c$ to
  $P_j$ alone, where $\featx^{A}_j=\featx$ is the public backbone and
  $\featx^{B}_j=\featx_j$ is the backbone composed with $P_j$'s adapter, and send
  $(\mathsf{Done},\mathsf{query},j)$ to $S$.
\end{enumerate}

\noindent\textbf{\emph{Leakage.}} $\Fhe$ reveals $\Leak$ to the adversary,
consisting of the public parameters, the per-client sample counts $\nj$, the
announced index $a^{\star}$, and the fact that each $\mathsf{Done}$ message was
sent.
\end{functionality}

The sample counts sit in the leakage because the protocol uses them as public
scalars. We prove security relative to both.

Three properties of $\Fhe$ are what the protocol must reproduce. No party ever
holds either shared head. The only value the whole federation learns is one index. A
client learns the labels it asks for, up to its allowance, and nothing else.

\subsection{Security Against Semi-Honest Parties}
\label{sec:semihonest}

\begin{definition}[Semi-honest security]
\label{def:semihonest}
Let $\mathcal{A}$ corrupt a set of parties that may contain the server and at
most $\tc-1$ clients. Write
$\mathsf{view}_{\mathcal{A}}^{\Pi}(\Dset)$ for the corrupted parties' joint view
of a protocol run on datasets $\Dset$. The protocol $\Pi$ securely realizes
$\Fhe$ with leakage $\Leak$ if there is a probabilistic polynomial-time
simulator $\mathcal{S}$ such that, for every $\Dset$,
\[
  \mathsf{view}_{\mathcal{A}}^{\Pi}(\Dset)
  \;\stackrel{c}{\approx}\;
  \mathcal{S}\big(\Leak,\ \{\Dj\}_{j\ \mathrm{corrupt}},\ \mathsf{out}_{\mathcal{A}}\big),
\]
where $\mathsf{out}_{\mathcal{A}}$ collects the answers $\Fhe$ returns to the
corrupted clients.
\end{definition}

\tronly{The simulator receives the answers because the protocol delivers them. A
definition that withheld them would be unachievable, and the theorem would say
nothing about a system that answers questions.}

\begin{theorem}[Semi-honest security]
\label{thm:semihonest}
Assume the $\tc$-out-of-$\Nc$ threshold CKKS scheme is secure under chosen
plaintext attack (IND-CPA) against an
adversary holding $\tc-1$ key shares. Then the protocol securely realizes $\Fhe$
with leakage $\Leak$ against any semi-honest adversary corrupting the server
and at most $\tc-1$ clients, in the sense of
\cref{def:semihonest}.
\end{theorem}

\begin{proof}[Proof sketch]
The simulator is given $\Leak$, the corrupted clients' datasets, and the answers
the corrupted clients receive.

\emph{Setup.} The simulator runs the distributed key generation with the
corrupted parties, playing the honest clients on uniformly chosen shares.


\emph{Training.} For each honest client the simulator sends an encryption of
zero in place of $\mathsf{Enc}(g_j\odot\disp)$, once for each of the two
arrangements, and in place of $\mathsf{Enc}(g_j)$, which the two share. The number
of ciphertexts follows from the public parameters, so the simulator knows it. A hybrid
argument replaces the honest ciphertexts one at a time, and each neighbouring
pair of hybrids is indistinguishable by IND-CPA against $\tc-1$ shares, the
assumption of the theorem.

\tronly{The server's additions and its multiplication by the reciprocal are
deterministic functions of those ciphertexts and of public scalars, so the
simulator computes them as the server does. The encrypted reciprocal is not such
a function. Its inversion circuit consumes levels and restores them by
collective refresh, an interactive protocol in which the participating clients
contribute fresh randomness. The simulator therefore invokes the simulator of
the refresh protocol for each of the two refreshes, playing the honest clients
on that protocol's simulated shares. Both refreshes sit at fixed points in the
circuit that do not depend on any plaintext, so the simulator knows their number
and position in advance.}

\paperonly{\emph{Selection and serving.} The selection step decrypts one value. The
simulator takes that index from $\Leak$.
For each query of a corrupted client, the simulator holds the answer in
$\mathsf{out}_{\mathcal{A}}$. Queries by honest clients are simulated as key switches of encryptions of zero.}

\tronly{\emph{Selection and serving.} The selection step decrypts one value. The
simulator takes that index from $\Leak$ and produces the collective
key-switching transcript for it from that protocol's simulator.
For each query of a corrupted client, the simulator holds the answer in
$\mathsf{out}_{\mathcal{A}}$ and simulates the key switch to that client in the
same way. Queries by honest clients are switched to keys the adversary does not
hold, so their transcripts are simulated as key switches of encryptions of zero.}

Composing the three parts gives a view computationally indistinguishable from
the real one. \paperonly{\trsee{sec:security} gives every proof in full.}
\end{proof}

\subsection{Input Privacy Against Malicious Clients}
\label{sec:malicious}

\Cref{thm:semihonest} assumes every party follows the protocol. We now let the
clients deviate. We bound what a deviating coalition learns about the data of a client
that does not deviate. We do not claim correctness.

\begin{definition}[Content game]
\label{def:contentgame}
The adversary $\mathcal{A}$ corrupts $\tc-1$ clients. The server is honest. At least one client $h$ stays honest.
\begin{enumerate}
\item $\mathcal{A}$ outputs two datasets $D_0$ and $D_1$ with $|D_0|=|D_1|$.
\item The challenger samples $b\in\{0,1\}$ and gives $D_b$ to client $h$.
\item The protocol runs. $\mathcal{A}$ may deviate arbitrarily in the values its
  clients upload, in the shares they contribute, and in the queries they ask.
\item $\mathcal{A}$ outputs $b'$. Its advantage is $|\Pr[b'=b]-\tfrac12|$.
\end{enumerate}
\end{definition}

\tronly{The equal-size restriction is necessary. The per-client sample counts are
public, so an adversary free to choose datasets of different sizes reads the
answer off the counts without attacking anything.}

\begin{theorem}[Input privacy]
\label{thm:malicious}
Assume the conditions of \cref{thm:semihonest}. Then the advantage of
$\mathcal{A}$ in \cref{def:contentgame} is at most
$\negl(\lambda)+\delta(Q_{\mathrm{tot}})$, where
$Q_{\mathrm{tot}}=(\tc-1)Q$ and $\delta$ is the advantage available from
$Q_{\mathrm{tot}}$ label queries to the served head.
\end{theorem}

\begin{proof}[Proof sketch]
The adversary's view consists of the key-generation transcript, the honest
clients' uploaded ciphertexts, the values the server computes, and the answers to its own queries.

\paperonly{The first four are independent of $b$ up to $\negl(\lambda)$.}

\tronly{The first four are independent of $b$ up to $\negl(\lambda)$. The honest
ciphertexts are indistinguishable from encryptions of zero by the hybrid
argument of \cref{thm:semihonest}, which holds against an adversary with
$\tc-1$ shares. The server is honest, so the values it computes are the
prescribed deterministic functions of those ciphertexts and of public scalars.
The server is honest, so every ciphertext presented for key switching is
the prescribed output of \cref{alg:serve} on a query, and its plaintext is a
label the coalition does not choose. Deviating in the uploaded values does not
help either, because the adversary already chooses those values in the real
protocol and they carry no dependence on $b$. Deviating in the contributed
shares does not help, because a share that is not the prescribed one makes the
reconstruction fail, which aborts the query and returns nothing.}

What remains is the answers. These do depend on $b$, because $\thstar$ depends
on client $h$'s displacement, and it is bounded by what $Q_{\mathrm{tot}}$ label queries
reveal.
\end{proof}

The bound scales with the corruption threshold. A coalition of $\tc-1$ clients
commands $(\tc-1)Q$ queries whatever the size of the federation.
\paperonly{\trsee{sec:experiments} measures $\delta$ and turns it into a bound on $Q$.}
\tronly{\Cref{sec:exp-leak} measures $\delta$ and turns it into a bound on $Q$.}

\subsection{Necessity of an Honest Server}
\label{sec:malicious-ext}

\Cref{thm:malicious} keeps the server honest. That assumption is necessary. Decryption is a protocol the clients run on a ciphertext the server presents. A
deviating server can present a ciphertext of its own choosing, such as a row of
$\thstar$, rather than the output of \cref{alg:serve}. The natural repair is to
have an honest client refuse such a request. It cannot, and the obstacle is the
encryption itself.

\begin{proposition}[No plaintext gate from a ciphertext]
\label{prop:nogate}
Let an honest client decide, by a polynomial-time predicate $D$ on its view,
whether to contribute its decryption share for a presented ciphertext. Let the
rest of that view be independent of the underlying plaintext. If the scheme is
IND-CPA against an adversary holding $\tc-1$ key shares, then for every pair of
plaintexts $m_0$ and $m_1$,
\[
  \big|\Pr[D(\mathsf{Enc}(m_0))=1]-\Pr[D(\mathsf{Enc}(m_1))=1]\big|
  \le \negl(\lambda).
\]
\end{proposition}

\begin{proof}
A gap between the two probabilities gives an IND-CPA distinguisher that holds
one key share. One share is at most $\tc-1$ shares, so the assumed security of
the scheme applies directly.
\end{proof}

An honest client therefore releases its share for both plaintexts or for
neither. No protocol of this message pattern realizes $\Fhe$ against a
malicious server, which is why \cref{thm:malicious} places the server among the
honest parties.

\paperonly{A mechanism checks provenance.
The serving circuit is a deterministic function of the encrypted head, the
encrypted query, the evaluation keys and the public parameters, and its level
restoration is deterministic as well. A
client can therefore recompute the prescribed output ciphertext and compare it as
a string.}

\tronly{Two mechanisms would remove the assumption, and we implement neither. The first
checks provenance rather than content. The serving circuit is a deterministic
function of the encrypted head, the encrypted query, the evaluation keys and the
public parameters, and its level restoration is deterministic as well, because
the protocol restores levels under collectively generated bootstrapping keys
rather than by collective refresh. A client can therefore recompute the
prescribed output ciphertext and compare it as a string. The design choice that
saves $510$\,MiB of traffic on every query is the same one that leaves the
circuit verifiable. Checking a random fraction $p$ of queries suffices, because
recovering the shared map needs on the order of $\Cc d$ deviating requests and a
campaign of that length escapes a single checker with probability at most
$(1-p)^{\Cc d}$, which at $p=0.01$ is below $10^{-13}$ on the smallest of our
tasks. The second mechanism certifies which function the server
evaluated, through a zero-knowledge proof of correct evaluation or verifiable
homomorphic encryption~\cite{viand2023verifiable}. \Cref{sec:scope} records both
as limitations.}

\subsection{Leakage Inherent to the Functionality}
\label{sec:inherent}

\tronly{\Cref{thm:semihonest} shows the protocol reveals no more than $\Fhe$ reveals.
That leaves the question of what $\Fhe$ itself reveals, and the answer is not
nothing.}

$\Fhe$ answers queries. Answers carry information about $\thstar$, and enough of
them reconstruct it. \Cref{sec:exp-leak} measures the rate at three to five
queries per parameter. This is a property of the
functionality. Any protocol that realizes $\Fhe$
inherits it.

The cryptographic claim is that the protocol adds nothing to the leakage of
$\Fhe$. The operational claim is that $\Fhe$'s own leakage is bounded by the query
allowance $Q$. The first rests on IND-CPA. The second rests on a measurement.

\tronly{\subsection{What a Fellow Client Learns}
\label{sec:c2c}}

\tronly{\Cref{thm:malicious} bounds a deviating coalition by
$\negl(\lambda)+\delta(Q_{\mathrm{tot}})$ and leaves $\delta$ unnamed. We name it
here. A coalition of clients is the one adversary the cryptography does not
silence, because its members hold key shares, run the backbone themselves, and
receive answers by design.}

\tronly{A client outside the coalition contributes four things to what the
coalition sees. It contributes the public sample count $\nj$, which the protocol
uses as a public scalar. It contributes the ciphertexts it uploads, which are
computationally independent of what they encrypt. It contributes its share of the
announced index $a^{\star}$ and of the phase signals. It contributes its part of
every answer the coalition is given. It contributes no per-class total, because
\cref{sec:split} inverts those under encryption and never decrypts them, and it
contributes no adapter, because no adapter aggregate exists.}

\tronly{The allowance alone does not bound what the answers reveal. Membership
inference from predicted labels is an established attack on a classifier that
answers queries~\cite{choquettechoo2021labelonly}, and its cost is not a function
the allowance controls. We therefore bound the channel at the head rather than at
the budget.}

\tronly{\begin{proposition}[The query channel is no stronger than the head]
\label{prop:c2c}
Let $\mathcal{A}$ corrupt at most $\tc-1$ clients in \cref{def:contentgame}, with
the server honest, and let it ask at most $Q_{\mathrm{tot}}=(\tc-1)Q$ queries.
Then there is an adversary $\mathcal{B}$ that receives both shared heads in
plaintext together with $\Leak$, asks no query, and satisfies
\[
  \mathsf{Adv}(\mathcal{A})\;\le\;\mathsf{Adv}(\mathcal{B})+\negl(\lambda).
\]
\end{proposition}}

\tronly{\begin{proof}
The proof of \cref{thm:malicious} splits the view of $\mathcal{A}$ into the
key-generation transcript, the honest clients' ciphertexts, the values the honest
server computes, and the answers to its own queries. The first three are
independent of $b$ up to $\negl(\lambda)$, and $\mathcal{B}$ produces them from
$\Leak$ and from encryptions of zero exactly as the simulator of
\cref{thm:semihonest} does.

Each answer is $\arg\max_c(\thstar_{a^{\star}}\,\featx^{a^{\star}}_j(x_i))_c$ for a
query $x_i$ that $\mathcal{A}$ asks at one of its own clients $j$. Under
$a^{\star}=A$ the map is the public backbone. Under $a^{\star}=B$ it composes the
backbone with the adapter of client $j$, and that client is corrupt, so
$\mathcal{B}$ holds the adapter. $\mathcal{B}$ holds both heads by assumption and
reads $a^{\star}$ from $\Leak$. It computes every answer itself.

$\mathcal{B}$ therefore runs $\mathcal{A}$, answers each query as the server would,
and returns whatever $\mathcal{A}$ returns. The simulation is exact outside the
negligible event on which the first three parts of the view differ.
\end{proof}}

\tronly{Write $\delta_{\mathrm{wb}}$ for the advantage of the best
polynomial-time adversary that holds both shared heads. \Cref{prop:c2c} gives
$\delta(Q_{\mathrm{tot}})\le\delta_{\mathrm{wb}}$ at every budget. What a
coalition learns about an honest client is capped by what the shared head itself
reveals, and that cap does not depend on the allowance. The allowance prices the
cost of approaching the cap, and \cref{sec:exp-leak} measures that cost.}

\tronly{Two consequences follow for a deployment. Lowering the threshold $\tc$
raises $Q_{\mathrm{tot}}$ and moves a coalition along the curve toward the cap,
and it does not move the cap. Raising the allowance does the same. Neither choice
changes what is ultimately learnable, so both are decisions about cost.}

\tronly{\subsubsection{Measuring the cap}
\label{sec:c2c-measured}}

\tronly{We measure $\delta_{\mathrm{wb}}$ and find it at chance. The adversary in
\cref{prop:c2c} receives both shared heads in plaintext, so we give an attacker
exactly that and run the likelihood-ratio membership attack of
Carlini et al.~\cite{carlini2022membership} against it. We build $64$ shadow
federations per cell, each one $\Nc$ clients training on half of their own
examples followed by the merge of \cref{eq:aggregate}, and we test $2000$
candidates over three tasks, three seeds and both arrangements. The attack
reaches an area under the curve between $0.49$ and $0.53$, and a true-positive
rate between $0.008$ and $0.023$ at a false-positive rate of $0.01$. A
true-positive rate equal to the false-positive rate is chance.}

\tronly{Handing the head over is a counterfactual, and that is the point of
measuring there. No party ever holds either head, so this is the world in which
the cryptography has already failed. \Cref{prop:c2c} says the query channel is no
stronger than that world. Measuring the strongest adversary the proposition
admits therefore bounds every adversary the protocol admits, and the bound is
chance.}

\tronly{A coalition that extracts a copy first does no better. We fit a copy of
the head to $2\times10^{5}$ label-only answers, the largest budget of
\cref{sec:exp-leak}, and attack the copy. The area under the curve moves to
between $0.49$ and $0.50$, and the true-positive rate at a false-positive rate of
$0.01$ to between $0.009$ and $0.017$. Extraction costs the attacker queries and
returns nothing, which is what \cref{prop:c2c} predicts and what no published
measurement had shown.}

\tronly{We report one figure we do not treat as a result. Predicting membership
from whether the head classifies an example correctly reaches an area under the
curve of $0.72$ on Banking77. Our holdout reserves at least one example of every
class a client holds, so held-out examples are richer in rare classes than
trained ones, and that rule separates them for a reason that has nothing to do
with membership. The likelihood-ratio attack calibrates each example against its
own shadow distributions and does not inherit the confound.}

\tronly{The coverage-weighted merge invites a sharper attack and does not yield
to it. A class that one client alone holds takes its row from that client's
displacement, with no averaging over the federation, and for a linear map with a
bias the ratio of the row displacement to the bias displacement is a weighted
mean of that client's features~\cite{phong2018privacy}. We compute that ratio for
all $570$ class rows of the three tasks. It matches the closest example of its own
class better than the closest example of another class by $0.012$ in cosine where
one client holds the class, and by $0.006$ where six or more do, against an
absolute cosine of $0.16$ and a random-direction baseline of $0.029$. The row
points into the space the features occupy. It does not point at a class, and it
does not point at a record.}

\tronly{Three other attack families ask for something this construction does not
provide, and we say which and why. Reconstruction from weights is the first. The
strongest published attack in our exact setting recovers training images from
heads trained on frozen $768$-dimensional embeddings, and its authors report that
it fails on linear probes and name linear models as a
mitigation~\cite{oz2024reconstructing}. Their own governing ratio is the
parameter count over the unknowns, which for a linear head reduces to the class
count over the client's sample count. Ours is about $0.01$, and the smallest
ratio at which they recover anything is $0.5$.}

\tronly{Attribute inference is the second. Against a model trained on features of
this kind, the black-box attack performs at the base rate, $0.28$ against a base
rate of $0.28$ on one dataset and $0.05$ against $0.06$ on
another~\cite{jayaraman2022attribute}. The variant that beats imputation reads
individual hidden-layer activations, and a single linear map has none.}

\tronly{Gradient inversion is the third, and it does not instantiate at all. Every
attack in that line needs gradients with respect to parameters that see the
input, or a model the server chose~\cite{zhu2019deep,geiping2020inverting}. The
backbone here is public and fixed, every client runs the same one, and the only
uploaded object is a trained displacement of a map that sits after the last
nonlinearity.}

\tronly{Two limits belong with these numbers. A thousand members resolve a
true-positive rate no finer than $0.001$, so every figure we report at a
false-positive rate of $0.001$ counts two to four examples rather than a rate.
And the shadow federations vary who trains on what while holding the coverage
pattern fixed, so they say nothing about how leakage moves when the partition
moves.}

\tronly{None of this makes the encryption idle, and the reading that it does
inverts the argument. The encryption is what stops any party from holding the
trained result, which \cref{sec:intro} gives as a requirement in its own right.
It is also what removes the per-round updates that carry the strongest published
attacks on federated learning~\cite{zhu2025fedmia}. What it leaves open, by
design, is a client asking for labels. These measurements close that last
channel, and they close it at the point where the cryptography is assumed to have
failed altogether.}

\tronly{\subsection{What a Decryption Carries}
\label{sec:cpad}}

\tronly{The protocol decrypts twice. The selection step decrypts one scalar, and
each answer is key-switched to the client that asked for it. Both are intended
outputs. In an exact scheme that would end the discussion, because a decryption
returns the message and nothing else. CKKS is approximate, so a decryption
returns the message together with an error term, and that term is a function of
the secret key and of the values the circuit consumed. Li and Micciancio show
that an adversary which observes decryption results recovers the key that
produced them~\cite{limicciancio2021approximate}.}

\tronly{Their attack does not need the message. It reads the ring element
$c_0+c_1 s$ out of the decoded output and inverts it, so an argument that the
adversary cannot predict what the circuit should have produced does not answer
it. The reason the attack does not reach us is structural. It recovers whichever
key performed the decryption. Every decryption here happens after a collective
key switch that moves the ciphertext off the collective key, so the key it
recovers is the querying client's own, which that client already holds. The
collective key $s$ never appears in a relation the coalition can invert, and
$\mathsf{Enc}(\thstar)$ is never decrypted at all.}

\tronly{The published threshold attack does not reach us either, for a separate
reason. It identifies ciphertexts whose noise is exactly zero and reads a linear
equation in the key off each one~\cite{checri2024threshold}. Detecting a
zero-noise ciphertext requires a rounding step whose success can be observed,
which exact schemes have and CKKS does not. Their statement about threshold CKKS
covers the case where no smudging is applied.}

\tronly{What remains is a condition rather than an attack, and we state it because
our parameters do not meet it. The key-switch simulator of the multiparty scheme
we build on is correct when the smudging noise dominates the noise of the
ciphertext being switched, which is what lets it produce an honest client's share
without knowing that noise~\cite{mouchet2021multiparty}. That scheme sets the
requirement at $2^{\lambda/2}$ times the ciphertext noise. Our implementation
smudges at eight times the fresh-encryption noise, which is the value the library
uses in its own tests, and the serving circuit leaves a ciphertext noise several
orders of magnitude larger than that. The simulation step of
\cref{thm:semihonest} therefore rests on a hypothesis our measured configuration
does not satisfy. What the step would otherwise hide is the circuit noise, which
depends on the honest clients' displacements. Whether that is exploitable is
open, and the scheme's authors record it as open as well.}

\tronly{Meeting the stated requirement is not a matter of raising one constant.
At a plaintext scale of $2^{45}$ the required smudging exceeds the scale itself,
so it would destroy the answer rather than protect it. A configuration that
satisfies the condition needs a larger scale and modulus, and every cost in
\cref{sec:exp-cost} would move with them. The tension is recorded in the
literature rather than particular to us. Threshold CKKS is insecure under this
notion when no smudging is applied, and flooding sized to the rule is expected
to cost the precision of the decrypted result~\cite{checri2024threshold}. We
report the configuration we measured and state the gap rather than claim a
guarantee we did not instantiate.}

\tronly{Repeating a query turns that gap into a key recovery, and we state the
arithmetic because it decides which control the deployment needs. Write the
served ciphertext as $(c_0,c_1)$ under the collective key $s$. The querier
decrypts $c_0+c_1s+E$, where $E$ collects the honest parties' smudging terms. The
serving circuit is deterministic, so a client that resubmits an identical query
receives the same $c_0$ and $c_1$ and a freshly drawn $E$. Averaging $k$
repetitions drives $E$ to zero and converges on $c_0+c_1s$ exactly. The noise of
the ciphertext is not a term outside that expression. It is inside it, because
$c_0+c_1s$ is the plaintext together with that noise. An adversary that holds
$c_0$ and $c_1$ subtracts and rounds, obtains $c_1s$, and recovers $s$ by one
ring inversion. The server holds $c_0$ and $c_1$, and \cref{def:semihonest}
corrupts the server together with $\tc-1$ clients, so this view is inside the
model rather than outside it. Resubmitting a query is also what the protocol
permits, so the adversary is semi-honest.}

\tronly{The cost is the number of repetitions that drives the averaged error
below half a unit in each of the $2^{15}$ coefficients, which is on the order of
$81\sigma^2$ for a resampled part of standard deviation $\sigma$. At eight times
the fresh-encryption noise and a quorum of ten that is about $5\times10^{5}$
repetitions, which is the same order as the query budgets \cref{sec:exp-leak}
already contemplates. A duplicate check is therefore a required control and not a
supplementary one. Re-randomizing a query defeats the check but not the control,
because a re-randomized query is a different ciphertext and yields one
independent sample rather than a repetition, and independent samples are the
regime the flooding is sized for. Sizing that flooding for $q$ decryptions costs
half a bit for each doubling of $q$, and the budget is refreshed by rekeying once
it is spent~\cite{limicciancio2022dp}.}

\tronly{Raising the smudging is the second control and it is free on every axis
the deployment measures. The values this protocol decrypts are one index and one
label, so the precision the answer needs is a handful of bits against a plaintext
scale of $2^{45}$. Moving the smudging from eight times the fresh-encryption
noise to $2^{25}$ multiplies the required repetitions by roughly $2^{40}$ and
still leaves twenty bits of the scale. It costs nothing in latency, nothing in
traffic, and nothing in the accuracy of the served answers, because the label is
an argmax and not a score. It does not meet the $2^{\lambda/2}$ rule, and we do
not claim that it does. It removes the attack above at no price, which is a
different and weaker statement, and it is the configuration we recommend.}

\tronly{One deviation the theorems do not cover belongs here as well. Collective
key generation has each party publish a share of the public key, and the shares
are summed. A party that publishes last, after seeing the others, can subtract
their contribution and leave a public key whose secret it holds alone. Every
ciphertext uploaded under that key is then readable by that party, and the
deviation is detected only later, when the honest decryption shares fail to
reconstruct, by which time the uploads have been read. The standard remedy is to
commit the shares before opening them, which costs one round and is independent
of the ring parameters. \Cref{thm:malicious} lets a coalition deviate in the
shares it contributes, and its proof disposes of that deviation for decryption
shares, where a wrong share aborts the query. Key-generation shares are not
covered by that argument, and a deployment that omits the commitment round is
outside the theorem.}
